% GENERATED FILE. Edit canonical lessons, then run npm run generate:books.
% canonical-source-hash: fnv1a64:6dc03e380152ced9
% canonical-lessons: FR-C123-infirmiere, FR-C123-cuisinier, FR-C123-policier, FR-C123-vendeur, FR-C123-agriculteur

\chapter{Nurse, Cook, Police officer, Shop assistant, Farmer}
\label{ch:fr-nurse-cook-police-officer-shop-assistant-farmer}
\hlchaptermodality{123}

\begin{chapteropening}
\textbf{By the end of this chapter:} \emph{I can name a nurse, a cook, a police officer, a shop assistant, a farmer.}

Complete the last lesson of chapter 123: i can name a nurse, a cook, a police officer, a shop assistant, a farmer.
\end{chapteropening}

\section[l'infirmière]{l'infirmière — a nurse}

\label{lesson:FR-C123-infirmiere}



\begin{quote}
Before the new one: say the French for a dessert, then the French for a meal.
\end{quote}

\subsection*{You'll want to know: l'infirmière}
\textbf{l'infirmière} — \textquotedblleft{}a nurse\textquotedblright{}.

\begin{grammarlens}[title={What you've built}]
One more word for this chapter.
\end{grammarlens}

\subsection*{Your turn}
\begin{itemize}
\raggedright
  \item \emph{Say it:} \emph{l'infirmière}
  \item \emph{Say it:} \emph{l'infirmière}, once more
  \item \emph{Say it:} say \emph{l'infirmière}
  \item \emph{Recall:} say the French for a dessert, then the French for a meal, then say \emph{l'infirmière} again
\end{itemize}

\begin{tcolorbox}[breakable,colback=teal!4,colframe=teal!35!black,title={Before you move on}]
What does l'infirmière mean? (\textquotedblleft{}A nurse\textquotedblright{}.) Say it once more.
\end{tcolorbox}
\section[le cuisinier]{le cuisinier — a cook}

\label{lesson:FR-C123-cuisinier}



\begin{quote}
Before the new one: say the French for a meal, then the French for a nurse.
\end{quote}

\subsection*{You'll want to know: le cuisinier}
\textbf{le cuisinier} — \textquotedblleft{}a cook\textquotedblright{}.

\begin{grammarlens}[title={What you've built}]
One more word for this chapter.
\end{grammarlens}

\subsection*{Your turn}
\begin{itemize}
\raggedright
  \item \emph{Say it:} \emph{le cuisinier}
  \item \emph{Say it:} \emph{le cuisinier}, once more
  \item \emph{Say it:} say \emph{le cuisinier}
  \item \emph{Recall:} say the French for a meal, then the French for a nurse, then say \emph{le cuisinier} again
\end{itemize}

\begin{tcolorbox}[breakable,colback=teal!4,colframe=teal!35!black,title={Before you move on}]
What does le cuisinier mean? (\textquotedblleft{}A cook\textquotedblright{}.) Say it once more.
\end{tcolorbox}
\section[le policier]{le policier — a police officer}

\label{lesson:FR-C123-policier}



\begin{quote}
Before the new one: say the French for a nurse, then the French for a cook.
\end{quote}

\subsection*{You'll want to know: le policier}
\textbf{le policier} — \textquotedblleft{}a police officer\textquotedblright{}.

\begin{grammarlens}[title={What you've built}]
One more word for this chapter.
\end{grammarlens}

\subsection*{Your turn}
\begin{itemize}
\raggedright
  \item \emph{Say it:} \emph{le policier}
  \item \emph{Say it:} \emph{le policier}, once more
  \item \emph{Say it:} say \emph{le policier}
  \item \emph{Recall:} say the French for a nurse, then the French for a cook, then say \emph{le policier} again
\end{itemize}

\begin{tcolorbox}[breakable,colback=teal!4,colframe=teal!35!black,title={Before you move on}]
What does le policier mean? (\textquotedblleft{}A police officer\textquotedblright{}.) Say it once more.
\end{tcolorbox}
\section[le vendeur]{le vendeur — a shop assistant}

\label{lesson:FR-C123-vendeur}



\begin{quote}
Before the new one: say the French for a cook, then the French for a police officer.
\end{quote}

\subsection*{You'll want to know: le vendeur}
\textbf{le vendeur} — \textquotedblleft{}a shop assistant\textquotedblright{}.

\begin{grammarlens}[title={What you've built}]
One more word for this chapter.
\end{grammarlens}

\subsection*{Your turn}
\begin{itemize}
\raggedright
  \item \emph{Say it:} \emph{le vendeur}
  \item \emph{Say it:} \emph{le vendeur}, once more
  \item \emph{Say it:} say \emph{le vendeur}
  \item \emph{Recall:} say the French for a cook, then the French for a police officer, then say \emph{le vendeur} again
\end{itemize}

\begin{tcolorbox}[breakable,colback=teal!4,colframe=teal!35!black,title={Before you move on}]
What does le vendeur mean? (\textquotedblleft{}A shop assistant\textquotedblright{}.) Say it once more.
\end{tcolorbox}
\section[l'agriculteur]{l'agriculteur — a farmer}

\label{lesson:FR-C123-agriculteur}



\begin{quote}
Before the new one: say the French for a police officer, then the French for a shop assistant.
\end{quote}

\subsection*{You'll want to know: l'agriculteur}
\textbf{l'agriculteur} — \textquotedblleft{}a farmer\textquotedblright{}.

\begin{grammarlens}[title={What you've built}]
One more word for this chapter.
\end{grammarlens}

\subsection*{Your turn}
\begin{itemize}
\raggedright
  \item \emph{Say it:} \emph{l'agriculteur}
  \item \emph{Say it:} \emph{l'agriculteur}, once more
  \item \emph{Say it:} say \emph{l'agriculteur}
  \item \emph{Recall:} say the French for a police officer, then the French for a shop assistant, then say \emph{l'agriculteur} again
\end{itemize}

\begin{tcolorbox}[breakable,colback=teal!4,colframe=teal!35!black,title={Before you move on}]
What does l'agriculteur mean? (\textquotedblleft{}A farmer\textquotedblright{}.) Say it once more.
\end{tcolorbox}
